\section{The Gram matrix of the product design}
\setcounter{equation}{0}

Recall $\bPsi(u) \in \R^{K}$, $K = 1+q\NJ$, the vector that collects the
constant and the wavelets of the $q$ modulators, and the identity
$\bZ = \Pi\{\bX \otimes \bPsi(\bU)\}$, where $\Pi$ is the fixed permutation
that moves the $p$ entries $X_{\cv}\cdot 1$ to the front. Since $\Pi$ is
orthogonal, every statement about eigenvalues is invariant to the column
order. Assumption 4 also gives a constant $C_\psi$, depending only on the
family, with
\begin{equation}
  \label{eq:pointwise}
  \sup_{u\in[0,1]}\sum_{\lev<J}\sum_{\tsl}\wav{\lev}{\tsl}(u)^2
  \;\le\; C_\psi\,2^{J} \qquad \text{for every } J \ge 1 :
\end{equation}
at each level at most $L$ periodized wavelets are nonzero at a given $u$ and
$\norm{\wav{\lev}{\tsl}}_{\infty} = 2^{\lev/2}\norm{\psi}_{\infty}$, so the
left side is at most
$L\norm{\psi}^2_{\infty}\sum_{\lev<J}2^{\lev} \le L\norm{\psi}^2_{\infty}2^J$.

\subsection*{Proof of Proposition 2}

\emph{Part (i), the modulators.} For $a \in \R^{K}$ write
$h_a(u) = a\T\bPsi(u) = a_0 + \sum_{\md}h_{a,\md}(u_{\md})$ with
$h_{a,\md} = \sum_{\lev\tsl}a_{\md,\lev\tsl}\wav{\lev}{\tsl}$. The functions
$1$ and $\wav{\lev}{\tsl}(u_{\md})$, $\md = 1,\dots,q$, are orthonormal in
$L_2([0,1]^q,du)$: within one coordinate this is the orthonormality of the
periodized system, and across coordinates
$\int\wav{\lev}{\tsl}(u_{\md})\wav{\lev'}{\tsl'}(u_{\md'})\,du =
\int_0^1\wav{\lev}{\tsl}\int_0^1\wav{\lev'}{\tsl'} = 0$ for
$\md \ne \md'$, because every wavelet integrates to zero. Hence
$\int_{[0,1]^q}h_a^2 = \norm{a}_2^2$ and, by Assumption 3,
$c_U\norm{a}_2^2 \le \E\{h_a(\bU)^2\} \le C_U\norm{a}_2^2$, which is the claim
since $a\T\Gram_{\bPsi}a = \E\{h_a(\bU)^2\}$.

\emph{Part (i), the full matrix.} Let $\delta \in \R^{pK}$ and reshape it as
$\Delta \in \R^{p\times K}$, row $\cv$ holding the coefficients attached to
$X_{\cv}$, so that $\{\bX\otimes\bPsi(\bU)\}\T\delta = \bX\T\Delta\bPsi(\bU)$.
By the tower property,
\[
  \delta\T\Pi\T\Gram\Pi\delta
  = \E\bigl[\{\bX\T\Delta\bPsi(\bU)\}^2\bigr]
  = \E\bigl[\bPsi(\bU)\T\Delta\T\,\E(\bX\bX\T\mid\bU)\,\Delta\bPsi(\bU)\bigr].
\]
Assumption 2 bounds the inner quadratic form between
$\kappa_1\norm{\Delta\bPsi(\bU)}_2^2$ and $\kappa_2\norm{\Delta\bPsi(\bU)}_2^2$
almost surely, and
$\E\norm{\Delta\bPsi(\bU)}_2^2 = \delta\T(I_p\otimes\Gram_{\bPsi})\delta$ lies
between $\lmin(\Gram_{\bPsi})\norm{\delta}_2^2$ and
$\lmaxx(\Gram_{\bPsi})\norm{\delta}_2^2$. Combining the two and using the
first part gives $\gamma = \kappa_1c_U \le \lmin(\Gram)$ and
$\lmaxx(\Gram) \le \kappa_2C_U$. If $\bX$ and $\bU$ are independent, the
conditional second moment matrix is constant, the same computation gives
$\delta\T\Pi\T\Gram\Pi\delta = \delta\T\{\E(\bX\bX\T)\otimes\Gram_{\bPsi}\}
\delta$ for every $\delta$, and two symmetric matrices with the same quadratic
form are equal; the eigenvalues of a Kronecker product of symmetric matrices
are the products of the eigenvalues, and
$\lmin\{\E(\bX\bX\T)\} \ge \kappa_1$ by Jensen's inequality, since $\lmin$ is
concave and $\E(\bX\bX\T) = \E\{\E(\bX\bX\T\mid\bU)\}$.

\emph{Part (ii), the empirical matrix.} First the range:
$\norm{\bZ_i}_2^2 = \norm{\bX_i}_2^2\norm{\bPsi(\bU_i)}_2^2 \le
pB_X^2(1+qC_\psi2^J) = R_J$ by \eqref{eq:pointwise} applied to each of the $q$
coordinates. Set $A_i = (\bZ_i\bZ_i\T-\Gram)/n$, independent, centered and
symmetric, with $\Gramhat - \Gram = \sum_i A_i$. Then
$\opnorm{\bZ_i\bZ_i\T} = \norm{\bZ_i}_2^2 \le R_J$ and
$\opnorm{\Gram} \le R_J$, so $\opnorm{A_i} \le 2R_J/n$; and
$\E\{(\bZ\bZ\T-\Gram)^2\} = \E(\norm{\bZ}_2^2\bZ\bZ\T) - \Gram^2 \preceq
R_J\Gram$, so that
$\opnorm{\sum_i\E A_i^2} \le R_J\lmaxx(\Gram)/n \le R_J\kappa_2C_U/n$. The
matrix Bernstein inequality \citep[Theorem 1.4]{tropp2012user} applied to
$\pm\sum_iA_i$ gives, for every $t>0$,
\[
  \Prob\bigl(\opnorm{\Gramhat-\Gram} \ge t\bigr)
  \;\le\; 2D\exp\Bigl\{-\frac{nt^2}{2R_J(\kappa_2C_U+2t/3)}\Bigr\},
\]
and $t = \gamma/2$ with Weyl's inequality gives
$\lmin(\Gramhat) \ge \lmin(\Gram)-\gamma/2 \ge \gamma/2$ on the complementary
event, which is the display of part (ii). The exponent is $-cn/R_J$ with $c$
fixed and $R_J \asymp pq2^J$, while $\log(2D) \asymp \log(pq2^J)$, so the
bound tends to zero exactly when $n/(pq2^J) - \log(pq2^J) \to \infty$, which
is implied by $pq2^J\log(pq2^J)/n \to 0$. On that event
$\compat^2(S;\Gramhat) \ge \gamma/2$ for every support, by
\eqref{eq:compat-lmin}. \qed

\begin{remark}[What the two assumptions exclude]
\label{rem:excl}
Assumption 2 fails exactly when a nontrivial linear combination of the linear
covariates is almost surely a function of the modulators on a set of positive
probability; in the canonical case $X_1\equiv1$ and $X_2 = h(U_1)$ the
conditional second moment matrix has rank one and the population Gram matrix
degenerates as $J$ grows, because $h(u_1)\wav{\lev}{\tsl}(u_1)$ is approached
by the span of the basis. In Assumption 3 it is the joint density that is
bounded below: if $U_2 = U_1$ almost surely, both marginals are uniform but
$\wav{\lev}{\tsl}(U_1)-\wav{\lev}{\tsl}(U_2) \equiv 0$ and
$\Gram_{\bPsi}$ is singular. Independence between $\bX$ and $\bU$ buys the
exact factorization of part (i) and nothing else.
\end{remark}

%%%%%%%%%%%%%%%%%%%%%%%%%%%%%%%%%%%%%%%%%%%%%%%%%%%%%%%%%%%%%%%%%%%%%%%%%%%%%%
% derivations/04-oraculo.tex
